\documentclass[11pt]{article}

\usepackage[margin=25mm]{geometry}

\usepackage{amsmath,amssymb,fontspec,unicode-math}

\IfFontExistsTF{DejaVu Serif}{\setmainfont{DejaVu Serif}}{\setmainfont{Latin Modern Roman}}

\IfFontExistsTF{Latin Modern Math}{\setmathfont{Latin Modern Math}}{\setmathfont{DejaVu Math TeX Gyre}}

\usepackage{xurl,hyperref}

\setlength{\parindent}{0pt}\setlength{\parskip}{6pt}

\begin{document}

\title{Spectral Analysis of Diffusion Models with Application to Schedule Design}

\author{ELVIA preliminary AI analysis}\date{21 October 2026}\maketitle

Preliminary assessment: No Issue Found.\par

Independent scientific review remains pending. Author adoption has not been requested.

No substantive error was established in the core shared-basis diffusion derivation. Finite KL optimization needs a nonsingularity/common-support qualification, and the eigenvalue-only Wasserstein expression needs the shared eigenbasis already present in that derivation; empirical recurrence/FID results were not refuted.

\section*{Review scope}

Rapid current-version reassessment of every substantive concern in the supplied LaTeX report: the implicated paper passages and relevant author-released source paths were inspected, existing proposed corrections were briefly checked, and impact was calibrated to the demonstrated scope. No author-provided or historical verification program was executed.

\section*{Problem 1: The finite KL schedule objective needs a nonsingularity or common-support condition}

Source: Section 3.3 (positive-semidefinite target covariance); Theorem 4.2, Eq. (16), PDF p. 6; Eq. (14) endpoint constraints.

Assessment: not\_issue; impact: no\_issue\_found

\subsection*{Source claim (paraphrase)}

Theorem 4.2, Eq. (16), PDF p. 6, gives a Gaussian KL schedule objective, while Section 3.3 permits positive-semidefinite target covariances. Appendix D.2, Eq. (38), PDF pp. 17–18, gives the usual determinant/inverse expression; the issue concerns its finite-domain qualification, not a main-text transcription sign.

\subsection*{Analysis}

A singular target Gaussian is supported on a proper affine subspace. If the generated Gaussian retains positive variance in a target-null direction because the final noise floor is nonzero, the measures do not have the required absolute continuity and KL is infinite. A pseudodeterminant alone cannot turn incompatible supports into a finite objective.

For \(\Lambda_0=\operatorname{diag}(0,1)\), mean \((1,0)\), and endpoints \(\bar\alpha_0=0.9999\), \(\bar\alpha_S=0.00004\), the null-direction DDIM variance telescopes to \((1-\bar\alpha_0)/(1-\bar\alpha_S)=5/49998>0\). The target assigns mass one to \(v_1=1\), while the generated Gaussian assigns zero there, so KL is \(+\infty\). This extended-real result does not contradict the recurrence; finite optimization needs positive-definite covariances or equal affine supports. The exact covariance and regularization used in the synthetic KL plot are unavailable, so a numerical error in that plot is not established.

\subsection*{Finding and scope}

Domain qualification for a finite optimization objective; no incorrect reported numerical covariance construction or core spectral recurrence was established.

\section*{Problem 2: The eigenvalue-only Wasserstein formula requires a shared eigenbasis}

Source: Appendix D.1, Eq. (36); common-basis Lemma 3.3.

Assessment: not\_issue; impact: no\_issue\_found

\subsection*{Source claim (paraphrase)}

Paraphrase: Appendix D.1 introduces the Wasserstein-\(2\) distance between two Gaussian distributions with means \(\mu_1,\mu_2\), covariances \(\Sigma_1,\Sigma_2\), and corresponding covariance eigenvalues using

\[W_2^2=\|\mu_1-\mu_2\|^2+ \sum_i\left(\sqrt{\lambda_i^{(1)}}-\sqrt{\lambda_i^{(2)}}\right)^2. \]

\subsection*{Analysis}

Covariances \(\operatorname{diag}(1,4)\) and \(\begin{pmatrix}2.5&-1.5\\-1.5&2.5\end{pmatrix}\) share eigenvalues but different orientations. Pairing only their sorted spectra would give zero, although their Gaussian Wasserstein distance is positive. The formula is valid under a shared eigenbasis, which the paper's central Gaussian derivation already supplies.

The unrestricted covariance calculation gives \(W_2^2=10-\sqrt{82}>0\), whereas sorted equal spectra give zero. In the paper’s own common eigenbasis, the Gaussian matrix-square-root formula reduces exactly to the displayed eigenvalue expression. The rotated example lies outside that already supplied common-basis condition.

\subsection*{Finding and scope}

Formula-scope qualification; the source core derivation already supplies its required common eigenbasis.

\section*{Limits}

No author-provided program, training procedure, model inference or wet-lab experiment was executed in this review.\par

Historical execution results are retained only with that attribution; a source-level defect alone does not identify the program or configuration used for the published experiments.\par

No author contact, adoption, independent human scientific review or novelty determination is established.\par

The findings apply only to the identified claims, source versions and stated assumptions; they do not certify or invalidate the entire paper.\par

\section*{Source identity}

arXiv:2502.00180v3 (9 December 2025), current author version

\url{https://arxiv.org/pdf/2502.00180v3}

Reviewed PDF SHA-256: \nolinkurl{902c7d4857c5bebd6aaff70d1bd4d733ab4b91105267ddace572a97a468f46b0}

The exact inspected current paper PDF is identified by its SHA-256. The official current arXiv history lists v3 dated 9 December 2025. Its exact PDF was read at the covariance-domain, KL-objective and shared-basis Wasserstein passages; both qualifications remain scoped rather than current paper errors. The current PDF spells the second author Michael Elad. An exact PDF acquisition timestamp is not established in this reassessment; source-currentness checks have their own recorded timestamps and file modification times were not substituted.

\end{document}
